\documentclass[11pt]{exam}
\input{../Header/packages.sty}
\input{../Header/styleExam.sty}
\input{../Header/mathexpressions.sty}
\input{../Header/headerExam.sty}
%=========================================================
\newcommand{\theme}{Trigonometrie}
\newcommand{\teacher}{Jonas Guler}
\newcommand{\class}{3Na}
\renewcommand{\date}{21.08.26}
\graphicspath{{Algebra/data-algebra},{Geometrie/data-geometrie},{Analysis/data-analysis}} % Setzt den Pfad zu den Bilder fest
%=========================================================
\begin{document}
\fontsize{14}{14}\selectfont
%========================================================
\noindent
\begingroup
\renewcommand{\arraystretch}{1.5}
\begin{tabular*}{\textwidth}{l @{\extracolsep{6pt}}l @{\extracolsep{\fill}}l @{\extracolsep{6pt}}l}
	\multicolumn{2}{l}{{\bf \theme}} & {\bf Note:} & {\answerbox}\\
	\multicolumn{2}{l}{{\bf \teacher}} & & \\
	{\bf Klasse:} & {\class} &  {\bf Name:} & {\rule{5cm}{0.5pt}}\\
	{\bf Dauer:} & $45$ Minuten & {\bf Datum:} & {\date}\\[8pt]
\end{tabular*}
\endgroup
\rule[1ex]{\textwidth}{1pt}
%=========================================================
\begin{center}
	\textbf{Bewertungstabelle}\\
	\small Wird von der Lehrperson ausgefüllt\\
	\vspace{5pt}
	\addpoints
	\gradetable[h][questions]\\
	\vspace{5pt}
	Der Rechenweg muss \textbf{vollständig und nachvollziehbar} sein.\\
	Lösen Sie die Prüfung mit einen schwarzen oder blauen Kugelschreiber.\\
	Als Hilfsmittel ist der Taschenrechner (TI - Nspire CX-II CAS) erlaubt.
	%Als Hilfsmittel ist der kleine Taschenrechner (TI-30X o.ä.) erlaubt.
\end{center}
\rule[1ex]{\textwidth}{1pt}

\vspace{10cm}
\begin{center}
	%\Huge{Repetitionsprüfung}
\end{center}

\newpage
\begin{questions}
	%=========================================================
	%---------------------TESTFRAGEN
	%=========================================================

	\newpage
	\question[5] %Geo2 S.45 A63
	Um die Breite eines Flusses (mit parallelen Ufern) zu berechnen, ist längs eines Ufers eine Strecke $\overline{AB}$ mit $50$m Länge bestimmt worden. Am gegenüberliegenden Ufer wird eine Stelle $P$ anvisiert, wobei $\alpha=64.1^\circ$ und $\beta=79.6^\circ$ messen. Berechne die Breite des Flusses.\\
	\emph{Tipp:} Mache eine sinnvolle Skizze der Situation.
	\vspace{7pt}

 \end{questions}
\end{document}